\begin{explanationbox}
	\textbf{\textcolor{CExplanation}{一句话直觉}}
	\textbf{通过绝对值使自变量非负，从而利用偶函数性质将非负部分的图像对称复制到负半轴。}

	\medskip
	\textbf{\textcolor{CExplanation}{核心拆解}}
	\begin{description}[style=nextline, leftmargin=2.8em, labelsep=0.5em, itemsep=0.45em]
		\item[\textbf{要点一（绝对值作用）：}]
			将自变量 $x$ 转化为 $|x|$，确保 $|x|\ge 0$，从而 $f(|x|)$ 在 $x<0$ 时实际使用正值 $-x$ 代入。
		\item[\textbf{要点二（偶函数判定）：}]
			对任意 $x$，$f(|-x|)=f(|x|)$，因此 $f(|x|)$ 是偶函数。
		\item[\textbf{要点三（图像生成规则）：}]
			在 $x\ge 0$ 时，$f(|x|)=f(x)$，图像保持不变；在 $x<0$ 时，$f(|x|)=f(-x)$，图像由 $x\ge 0$ 部分关于 $y$ 轴对称得到。
	\end{description}

	\medskip
	\textbf{\textcolor{CExplanation}{几何本质（左右对称性）}}
	\textbf{函数的右半平面图像完全决定整体图像；变换后的图像关于 $y$ 轴对称。}

	\medskip
	\textbf{\textcolor{CExplanation}{代数意义（函数值拼合）}}
	\textbf{$f(|x|)$ 由 $f(x)$ 和 $f(-x)$ 拼合而成，但 $f(-x)$ 又由 $f(x)$ 在正半轴的值决定，因此本质上是 $f(x)$ 在 $[0,+\infty)$ 上的值按偶函数规则拓展到全体实数。}

	\medskip
	\textbf{\textcolor{CExplanation}{考点价值（命题角度）}}
	\begin{itemize}[leftmargin=*, itemsep=0.5em, topsep=0.4em]
		\item \textbf{考法一（表达式转化）：} 已知 $f(x)$ 在 $x\ge 0$ 的解析式，求 $f(|x|)$ 的分段表示。
		\item \textbf{考法二（图像应用）：} 利用翻折规则快速画图，判断值域、单调性等。
		\item \textbf{考法三（奇偶性判定）：} 用 $f(|x|)$ 的偶函数性质证明或判定函数的奇偶性。
	\end{itemize}

	\medskip
	\textbf{\textcolor{CExplanation}{顿悟点}}
	\textbf{绝对值是对自变量做“非负化”处理，进而自动为函数赋予偶对称性。}

	\medskip
	\textbf{\textcolor{CExplanation}{使用场景}}
	\begin{itemize}[leftmargin=*, itemsep=0.5em, topsep=0.4em]
		\item \textbf{场景一（求解析式）：} 给定 $f(x)$ 在 $x\ge 0$ 的表达式，写出 $f(|x|)$ 的解析式。
		\item \textbf{场景二（图像对称）：} 已知 $f(x)$ 在 $x\ge 0$ 的图像，直接绘制 $f(|x|)$ 的图像（保留右侧，翻折左侧）。
		\item \textbf{场景三（单调性分析）：} 利用 $f(|x|)$ 在对称区间上的单调性相反的性质，快速判断单调区间。
	\end{itemize}

\end{explanationbox}
